\begin{examplebox}
	\textbf{\textcolor{CExample}{例1（基础）}}
	\textbf{\textcolor{CExample}{题目：}}
	在棱长为$1$的正方体$ABCD-A_1B_1C_1D_1$中，求异面直线$A_1B$与$B_1C$所成角的余弦值。
	\textbf{\textcolor{CExample}{解题步骤：}}
	\begin{description}[style=nextline, leftmargin=2.8em, labelsep=0.5em, itemsep=0.45em]
		\item[\textbf{第一步（平移构造）：}]
			平移$A_1B$至$D_1C$，则$D_1C$与$B_1C$交于点$C$，记$m=D_1C$，$n=B_1C$，$d=D_1B_1$。
			\par\textbf{理由：} 平移不改变线段长度和方向，$A_1B\parallel D_1C$，故异面直线所成角等于$D_1C$与$B_1C$的夹角。
		\item[\textbf{第二步（长度计算）：}]
			计算$m,n,d$的长度。因为棱长为$1$，所以$m=D_1C=\sqrt{2}$，$n=B_1C=\sqrt{2}$，$d=D_1B_1=\sqrt{2}$。
			\par\textbf{理由：} 在正方体中，面对角线长度均为$\sqrt{2}$。
		\item[\textbf{第三步（套用公式）：}]
			代入公式$\cos\theta=\frac{|m^2+n^2-d^2|}{2mn}$，得
			\[
				\cos\theta=\frac{|(\sqrt{2})^2+(\sqrt{2})^2-(\sqrt{2})^2|}{2\cdot\sqrt{2}\cdot\sqrt{2}}=\frac{2+2-2}{2\cdot2}=\frac{2}{4}=\frac12.
			\]
			所以$\cos\theta=\frac12$，$\theta=60^\circ$。
			\par\textbf{理由：} 直接代入数据计算并化简。
	\end{description}
	\textbf{\textcolor{CExample}{关键结论：}}
	\textbf{$\cos\theta=\frac12$，异面直线所成角为$60^\circ$。}

	\medskip
	\textbf{\textcolor{CExample}{例2（稍变形）}}
	\textbf{\textcolor{CExample}{题目：}}
	在正四棱柱$ABCD-A_1B_1C_1D_1$中，$AB=2$，$AA_1=1$，求异面直线$A_1B$与$B_1D_1$所成角的余弦值。
	\textbf{\textcolor{CExample}{解题步骤：}}
	\begin{description}[style=nextline, leftmargin=2.8em, labelsep=0.5em, itemsep=0.45em]
		\item[\textbf{第一步（平移构造）：}]
			平移$A_1B$至$D_1C$，则$D_1C$与$B_1D_1$交于点$D_1$，记$m=D_1C$，$n=B_1D_1$，$d=B_1C$。
			\par\textbf{理由：} $A_1B\parallel D_1C$，故异面直线所成角等于$D_1C$与$B_1D_1$的夹角。
		\item[\textbf{第二步（长度计算）：}]
			计算三边长度。在正四棱柱中，底面正方形边长为$2$，侧棱长为$1$，则 $m=D_1C=A_1B=\sqrt{AB^2+AA_1^2}=\sqrt{2^2+1^2}=\sqrt{5}$， $n=B_1D_1=\sqrt{AB^2+BC^2}=\sqrt{2^2+2^2}=2\sqrt{2}$， $d=B_1C=\sqrt{BC^2+CC_1^2}=\sqrt{2^2+1^2}=\sqrt{5}$。
			\par\textbf{理由：} 利用矩形对角线公式计算各线段长度。
		\item[\textbf{第三步（套用公式）：}]
			代入公式$\cos\theta=\frac{|m^2+n^2-d^2|}{2mn}$，得
			\[
				\cos\theta=\frac{|(\sqrt{5})^2+(2\sqrt{2})^2-(\sqrt{5})^2|}{2\cdot\sqrt{5}\cdot2\sqrt{2}}=\frac{|5+8-5|}{4\sqrt{10}}=\frac{8}{4\sqrt{10}}=\frac{2}{\sqrt{10}}=\frac{\sqrt{10}}{5}.
			\]
			所以异面直线所成角的余弦值为$\frac{\sqrt{10}}{5}$。
			\par\textbf{理由：} 代入数值计算并化简。
	\end{description}
	\textbf{\textcolor{CExample}{关键结论：}}
	\textbf{$\cos\theta=\frac{\sqrt{10}}{5}$。}

	\medskip
	\textbf{\textcolor{CExample}{例3（综合应用）}}
	\textbf{\textcolor{CExample}{题目：}}
	在直三棱柱$ABC-A_1B_1C_1$中，$\angle ABC=90^\circ$，$AB=BC=1$，$AA_1=2$，求异面直线$A_1B$与$AC_1$所成角的余弦值。
	\textbf{\textcolor{CExample}{解题步骤：}}
	\begin{description}[style=nextline, leftmargin=2.8em, labelsep=0.5em, itemsep=0.45em]
		\item[\textbf{第一步（平移转化）：}]
			连接$AB_1$，则$A_1B\parallel AB_1$，所以异面直线$A_1B$与$AC_1$所成角等于$AB_1$与$AC_1$的夹角。记$m=AB_1$，$n=AC_1$，$d=B_1C_1$。
			\par\textbf{理由：} 在直三棱柱中，侧面$ABB_1A_1$为矩形，故$A_1B\parallel AB_1$。
		\item[\textbf{第二步（长度计算）：}]
			由已知，$AB=BC=1$，$AA_1=2$，$\angle ABC=90^\circ$，则 $m=AB_1=\sqrt{AB^2+BB_1^2}=\sqrt{1^2+2^2}=\sqrt{5}$， $n=AC_1=\sqrt{AC^2+CC_1^2}=\sqrt{(\sqrt{2})^2+2^2}=\sqrt{2+4}=\sqrt{6}$， $d=B_1C_1=BC=1$。
			\par\textbf{理由：} 利用矩形和直角三角形性质计算各线段长度。
		\item[\textbf{第三步（套用公式）：}]
			在$\triangle AB_1C_1$中，由公式得
			\[
				\cos\theta=\frac{|m^2+n^2-d^2|}{2mn}=\frac{|5+6-1|}{2\cdot\sqrt{5}\cdot\sqrt{6}}=\frac{10}{2\sqrt{30}}=\frac{5}{\sqrt{30}}=\frac{\sqrt{30}}{6}.
			\]
			因此异面直线$A_1B$与$AC_1$所成角的余弦值为$\frac{\sqrt{30}}{6}$。
			\par\textbf{理由：} 代入$m,n,d$的数值，计算并化简。
	\end{description}
	\textbf{\textcolor{CExample}{关键结论：}}
	\textbf{$\cos\theta=\frac{\sqrt{30}}{6}$。}
\end{examplebox}
